% Consolidated 34-atomic-scattering from 620defd2149c8dcfd52d29c147aa4708c27ab8f3; original TeX fragments preserved.

% === BEGIN 34-atomic-scattering-001-018.tex blob a5f8ec83315e1bc30fcc6b25bfb8ff674f6c1cee ===
% Irodov Part VI §6.1 problems 1--18.
% Source: user-provided Arihant PDF pp.241--243 (printed pp.232--234),
% Answers pp.348--349, numerical input verified against page scans.
% Equations here use SI; C=Z1 Z2 e^2/(4 pi eps0) converts Gaussian e^2.
\atomicproblem{1}{232}
用汤姆孙的原子模型估算氢原子半径 $R$，
以及电子振动发出的光波长。已知电离能
$E_{\mathrm{ion}}=13.6\,\mathrm{eV}$。
\exsolution
汤姆孙模型将 $+e$ 均匀分布在半径 $R$ 的球内，
电子在球心稳定。电子由球心移至无穷远的
势能增加为
$E_{\mathrm{ion}}=3e^2/(8\pi\varepsilon_0R)$，
故
\[
 \boxed{R=\frac{3e^2}{8\pi\varepsilon_0E_{\mathrm{ion}}}
 \simeq0.159\,\mathrm{nm}}.
\]
在均匀正电球内电场 $E(r)=er/(4\pi\varepsilon_0R^3)$，
电子受到的回复力 $F=-e^2r/(4\pi\varepsilon_0R^3)$，
所以
\[
 \omega=\sqrt{\frac{e^2}{4\pi\varepsilon_0m_eR^3}},
 \quad\boxed{\lambda=\frac{2\pi c}{\omega}
 \simeq0.24\,\mu\mathrm m}.
\]
\excheck 本题严格依\emph{经典汤姆孙模型}计算，
不应拿现代量子氢原子的 Bohr 半径替代 $R$。

\atomicproblem{2}{232}
动能 $T=0.27\,\mathrm{MeV}$ 的 $\alpha$ 粒子
在金箔中被偏转 $\theta=60^\circ$。
求对应的碰撞瞄准参数 $b$。
\exsolution
对固定重核的库仑散射采用卢瑟福关系
\[
 \tan\frac\theta2=
 \frac{C}{2bT},\qquad
 C=\frac{(2e)(79e)}{4\pi\varepsilon_0}.
\]
故
\[
 \boxed{b=\frac C{2T}\cot\frac\theta2}.
\]
取 $e^2/(4\pi\varepsilon_0)
=1.440\,\mathrm{MeV\,fm}$ 得
\[
 \boxed{b\simeq7.30\times10^2\,\mathrm{fm}
 =0.730\,\mathrm{pm}}.
\]
\excheck 底本第六篇采用高斯单位制；
本译稿用 SI，故 $C$ 中明确包含 $1/(4\pi\varepsilon_0)$。

\atomicproblem{3}{232--233}
动能 $T=0.40\,\mathrm{MeV}$ 的 $\alpha$ 粒子
正面撞向 (a) 固定铅核；(b) 起初静止的\emph{自由}
$^{7}\mathrm{Li}$ 核。分别求最近距离。
\exsolution
设库仑能常数
$C=(2Z)e^2/(4\pi\varepsilon_0)$。
(a) 固定核没有反冲，初始动能在最近点全转为静电势能，
$T=C/r_{\min}$，其中 $Z_{\mathrm{Pb}}=82$。
因此
\[
 \boxed{r_{\min,\mathrm{Pb}}=\frac{2Z_{\mathrm{Pb}}k_e e^2}{T}
 \simeq0.590\,\mathrm{pm}}.
\]
(b) 自由锂核也会反冲，质心系相对动能为
$T_{\rm rel}=T\,M_{\rm Li}/(M_{\rm Li}+m_\alpha)$。
因此
\[
 \boxed{r_{\min,\mathrm{Li}}=
 \frac{2Z_{\mathrm{Li}}k_e e^2}{T}
 \left(1+\frac{m_\alpha}{M_{\mathrm{Li}}}\right)
 \simeq0.034\,\mathrm{pm}},
 \quad Z_{\mathrm{Li}}=3.
\]
\excheck $k_e=1/(4\pi\varepsilon_0)$；
自由靶核的修正因子大于1，不能直接用静止不动核的公式。

\atomicproblem{4}{233}
动能 $T=0.50\,\mathrm{MeV}$ 的 $\alpha$ 粒子
在固定汞核（$Z=80$）的库仑场中偏转
$\theta=90^\circ$。求轨道最小曲率半径和最近距离。
\exsolution
仍记 $C=2Zk_e e^2$。由库仑双曲线轨道几何
及卢瑟福散射公式可得
\[
 \boxed{\mathcal R_{\min}=
 \frac C{2T}\cot^2\frac\theta2},\qquad
 \boxed{r_{\min}=
 \frac C{2T}\left(1+\csc\frac\theta2\right)}.
\]
这里 $C=230.4\,\mathrm{MeV\,fm}$，
$\cot45^\circ=1$，
\[
 \boxed{\mathcal R_{\min}\simeq0.230\,\mathrm{pm}},
 \qquad\boxed{r_{\min}\simeq0.556\,\mathrm{pm}}.
\]
\excheck 两种长度不相同；
最近曲率半径属于轨迹局部几何，最近距离是到核的径向距离。

\atomicproblem{5}{233}
动能为 $T$、瞄准参数为 $b$ 的质子
被起初静止的金核散射，求传给金核的动量大小。
采用重核近似。
\exsolution
入射质子动量为 $p=\sqrt{2m_pT}$，
散射前后速率在固定靶近似下不变，
动量变化量为
$|\Delta\boldsymbol p|=2p\sin(\theta/2)$。
由
$\tan(\theta/2)=C/(2bT)$、
$C=79k_e e^2$ 得
\[
 \sin\frac\theta2=
 \frac1{\sqrt{1+(2bT/C)^2}}.
\]
根据动量守恒，金核所获动量大小为
\[
 \boxed{P_{\rm Au}=
 \frac{2\sqrt{2m_pT}}{\sqrt{1+(2bT/C)^2}}},
 \qquad C=79k_e e^2.
\]
\excheck 核的反冲动量方向与质子动量变化反向；
若求精细反冲能量，需要采用两体约化质量修正。

\atomicproblem{6}{233}
动能 $T=10\,\mathrm{MeV}$ 的质子以直线轨迹
掠过静止自由电子，最近距离 $b=10\,\mathrm{pm}$。
假定质子路径不偏折、电子在通过过程中几乎不动。
求电子得到的能量。
\exsolution
令质子速度
$v=\sqrt{2T/m_p}$。
电子的横向力为
$F_\perp(t)=k_ee^2 b/(b^2+v^2t^2)^{3/2}$。
沿完整经过时间积分，
\[
 \Delta p_\perp=\int_{-\infty}^{+\infty}F_\perp(t)\,dt
 =\frac{2k_ee^2}{bv}.
\]
电子获得的动能
\[
 T_e=\frac{(\Delta p_\perp)^2}{2m_e}
 =\boxed{\frac{m_p(k_ee^2)^2}{m_eb^2T}
 \simeq4.0\,\mathrm{eV}}.
\]
\excheck 纵向冲量为奇函数积分，等于0；
此题是小反冲、直线冲量近似而非完整双体散射。

\atomicproblem{7}{233}
粒子动能为 $T$、瞄准参数为 $b$，
穿过半径 $R$、深度 $U_0$ 的球形方势阱：
$U=0$（$r>R$）、$U=-U_0$（$r<R$）。
求偏转角 $\theta$ 与 $b$ 的关系。
\exsolution
粒子在阱内的速度比阱外大，
利用轨道的切向动量（角动量）守恒，
等效折射率
$n=v_{\rm in}/v_{\rm out}
=\sqrt{1+U_0/T}>1$。
令入射线与界面法线的夹角为 $i$，
内部为 $r$，有
$\sin i=b/R$、$\sin r=b/(nR)$。
两次穿越界面使轨迹总偏折
\[
 \boxed{\theta=
 2\left[\arcsin\frac bR-
 \arcsin\frac b{nR}\right]},
 \qquad 0\le b\le R.
\]
反解
\[
 \boxed{b=\frac{nR\sin(\theta/2)}
 {\sqrt{1+n^2-2n\cos(\theta/2)}}}.
\]
若 $b>R$ 则不进入阱，偏转角为0。
\excheck $U_0\to0$ 时 $n\to1$，
$\theta\to0$，符合无势阱极限。

\atomicproblem{8}{233}
半径 $R$ 的固定硬球受到半径为 $r$ 的
均匀平行粒子流照射，粒子与球作理想弹性碰撞。
求：(a) 偏角 $\theta$ 与瞄准参数 $b$；
(b) 碰后粒子落在 $(\theta,\theta+d\theta)$ 的条件概率；
(c) 偏角小于 $\pi/2$ 的概率。
\exsolution
把有限尺寸的入射粒子与球等效为
有效半径 $a=R+r$ 的点粒子镜面碰撞。
几何法线使碰前碰后速度对称，故
\[
 \boxed{\cos\frac\theta2=\frac ba},\quad 0\le b\le a.
\]
对于均匀入射、已发生碰撞的粒子，
撞击参数分布是
$dP=2b\,db/a^2$。
由 $b=a\cos(\theta/2)$ 微分并取绝对值，
\[
 \boxed{dP=\frac12\sin\theta\,d\theta}.
\]
积分 $0\le\theta<\pi/2$，
\[
 \boxed{P(\theta<\pi/2)=
 \int_0^{\pi/2}\tfrac12\sin\theta\,d\theta=\frac12}.
\]
\excheck (b)(c) 的概率以\emph{发生碰撞}为条件，
不涉及整个无限截面入射束中未碰撞的粒子。

\atomicproblem{9}{233}
动能 $T=1.0\,\mathrm{MeV}$ 的窄 $\alpha$ 束
正入射厚 $d=1.0\,\mu\mathrm m$ 的铂箔。
一个探测计数器位于偏转角 $60^\circ$、距靶
$\ell=10\,\mathrm{cm}$ 处，入口面积
$S=1.0\,\mathrm{cm^2}$。
求总入射粒子中进入计数器的比例。
\exsolution
铂的原子序数 $Z=78$、
密度 $\rho=21.5\,\mathrm{g/cm^3}$、
摩尔质量 $M\simeq195\,\mathrm{g/mol}$。
单位面积靶核数
$n_A=\rho dN_A/M$。
在单次散射、薄靶和小探测立体角近似下
\[
 \frac{\Delta N}{N}=
 n_A\left(\frac C{4T}\right)^2
 \csc^4\frac\theta2\,\Delta\Omega,\qquad
 C=2Zk_ee^2,\quad \Delta\Omega\simeq\frac S{\ell^2}.
\]
代 $\Delta\Omega=0.010\,\mathrm{sr}$ 得
\[
 \boxed{\frac{\Delta N}{N}\simeq3.3\times10^{-5}}.
\]
\excheck 必须将面积、长度和靶材面核密度统一单位；
此为\emph{探测几何}内的概率而不是总散射概率。

\atomicproblem{10}{233--234}
动能 $T=0.50\,\mathrm{MeV}$、通量
$\dot N_0=5.0\times10^3\,\mathrm{s^{-1}}$
的窄 $\alpha$ 束正入射金箔。
在 $\theta=60^\circ$、距靶 $\ell=15\,\mathrm{cm}$
处，测得粒子流密度
$J=40\,\mathrm{cm^{-2}s^{-1}}$。
求金箔厚度 $d$。
\exsolution
靶核体积数密度
$n_V=\rho_{\rm Au}N_A/M_{\rm Au}$，
单次卢瑟福散射预测的计数器位置通量为
\[
 J=\frac{\dot N_0}{\ell^2}
 n_Vd\left(\frac C{4T}\right)^2
 \csc^4\frac\theta2,\quad C=2(79)k_ee^2.
\]
故
\[
 \boxed{d=\frac{J\ell^2}{\dot N_0n_V}
 \left(\frac{4T}{C}\right)^2
 \sin^4\frac\theta2}.
\]
直接代入\emph{本版印刷题面}
$\rho_{\rm Au}\simeq19.3\,\mathrm{g/cm^3}$ 得
\[
 \boxed{d\simeq1.47\times10^{-4}\,\mathrm m
 =147\,\mu\mathrm m}.
\]
\excheck 本版书末印\emph{约 $1.5\,\mu\mathrm m$}，
与题面按上述薄靶公式计算相差\emph{约100倍}；
题面 $\dot N_0$ 与 $J$ 已按扫描页复查，暂保留该差异。
由此反推的厚箔已可能违反“单次散射薄靶”前提，
因此该数值不得标记为完全物理自洽。

\atomicproblem{11}{234}
一束 $\alpha$ 粒子依次照射质量面密度相同的
银箔和铂箔，后者相同角区计数率为前者的
$\eta=1.52$ 倍。
求铂原子序数（已知银原子序数与两材料原子量）。
\exsolution
相同质量面密度时核的面数密度与摩尔质量
$M$ 成反比；同样束流、能量与角度下
卢瑟福散射微分截面与 $Z^2$ 成正比。
因而
\[
 \eta=\frac{Z_{\rm Pt}^2/M_{\rm Pt}}
 {Z_{\rm Ag}^2/M_{\rm Ag}},\quad
 \boxed{Z_{\rm Pt}=Z_{\rm Ag}
 \sqrt{\eta\,\frac{M_{\rm Pt}}{M_{\rm Ag}}}}.
\]
以 $Z_{\rm Ag}=47$、
$M_{\rm Pt}\simeq195$、
$M_{\rm Ag}\simeq108$ 代入，
\[
 \boxed{Z_{\rm Pt}\simeq78}.
\]
\excheck 本版原题为 $\eta=1.52$，
不是 PDF 文本层中容易误读的“152”。

\atomicproblem{12}{234}
$\alpha$ 粒子束动能
$T=0.50\,\mathrm{MeV}$、
入射率 $I_0=5.0\times10^6\,\mathrm{s^{-1}}$，
正入射质量面密度
$\rho d=1.5\,\mathrm{mg/cm^2}$ 的金箔，
持续 $\tau=30\,\mathrm{min}$。
求散射角落在(a) $59^\circ$--$61^\circ$；
(b) 大于 $60^\circ$ 的粒子数。
\exsolution
靶核面数密度为
$n_A=\rho dN_A/M_{\rm Au}$，
总入射数 $N_0=I_0\tau$。
卢瑟福微分截面
$d\sigma/d\Omega=(C/(4T))^2\csc^4(\theta/2)$，
$C=2(79)k_ee^2$。
在极角区间 $\theta_a\le\theta\le\theta_b$
对全方位积分，得
\[
 \sigma_{ab}=\pi\left(\frac C{2T}\right)^2
 \left[\cot^2\frac{\theta_a}{2}
 -\cot^2\frac{\theta_b}{2}\right].
\]
因此
\[
 \boxed{N_{ab}=I_0\tau n_A\sigma_{ab}}.
\]
按\emph{扫描原题数值}代入，
\[
 \boxed{N_{59^\circ\text{--}61^\circ}
 \simeq1.62\times10^7},\qquad
 \boxed{N_{\theta>60^\circ}\simeq2.01\times10^8}.
\]
\excheck 本版书末分别列约 $1.6\times10^6$、
$2.0\times10^7$，低于照印刷题面算出的结果10倍。
原书扫描已确认质量面密度为 $1.5\,\mathrm{mg/cm^2}$，
故暂按原题算式并登记书末冲突，不自行把参数改小10倍。

\atomicproblem{13}{234}
速度 $v=6\times10^6\,\mathrm{m/s}$ 的窄质子束，
正入射厚 $d=1.0\,\mu\mathrm m$ 的银箔。
求后向半球（散射角 $\theta>90^\circ$）
的散射概率。
\exsolution
质子动能
$T=\tfrac12m_pv^2\simeq0.188\,\mathrm{MeV}$。
取银 $Z=47$，
$C=Zk_ee^2$，
靶核面数密度
$n_A=\rho_{\rm Ag}dN_A/M_{\rm Ag}$。
对 $\theta>90^\circ$ 积分卢瑟福截面：
\[
 \sigma_{\mathrm{back}}=
 \pi\left(\frac C{2T}\right)^2\cot^2(45^\circ)
 =\pi\left(\frac C{2T}\right)^2.
\]
故
\[
 \boxed{P=n_A\sigma_{\mathrm{back}}
 \simeq6.0\times10^{-3}}.
\]
\excheck 概率小于1，单次散射近似在此角区合理。

\atomicproblem{14}{234}
动能 $T=600\,\mathrm{keV}$ 的 $\alpha$ 束
正入射金箔，面核密度为
$n_A=1.1\times10^{19}\,\mathrm{cm^{-2}}$。
求偏转角小于 $\theta_0=20^\circ$
的束流份额（包括几乎不被散射的直行粒子）。
\exsolution
超过阈角的单次散射截面为
\[
 \sigma_{>\theta_0}=
 \pi\left(\frac C{2T}\right)^2
 \cot^2\frac{\theta_0}{2},\quad C=2(79)k_ee^2.
\]
薄靶近似下偏角超过阈值的粒子份额
$P_>=n_A\sigma_{>\theta_0}$。
因此偏角\emph{小于}阈值的份额为
\[
 \boxed{P_<\simeq1-n_A\pi\left(\frac C{2T}\right)^2
 \cot^2\frac{\theta_0}{2}\simeq0.60}.
\]
\excheck 单纯从$\theta=0$积分卢瑟福公式会发散，
应统计“没有发生大角散射”的补事件。
\emph{本题的一阶薄靶近似尤其需要误差警示：}
按照本题参数，大角单次散射的光学厚度
$\tau_{>20^\circ}=n_A\sigma_{>20^\circ}\simeq0.40$，
并非远小于1。
所以$1-\tau\simeq0.60$只能看作
\emph{一阶近似}。若仅考虑多次独立
“超过20°的事件”并做泊松重求和，
零此类事件的概率是$e^{-\tau}\simeq0.67$，
两种计算已有约7个百分点差距；
真实多重小角散射还可能合成超过20°的总偏角。
因此\emph{不能}把0.60或0.67
都误称为无需考虑多重散射的严格束流概率。
准确角分布必须给定更完整的多重散射模型。

\atomicproblem{15}{234}
动能 $T=1.4\,\mathrm{MeV}$ 的质子束照射黄铜箔，
质量面密度 $\rho d=1.5\,\mathrm{mg/cm^2}$，
铜与锌质量比为 $7:3$。
求偏角超过 $\theta_0=30^\circ$ 的粒子份额。
\exsolution
对元素 $j=\mathrm{Cu},\mathrm{Zn}$，阈角以上截面
\[
 \sigma_j(>\theta_0)=
 \pi\left(\frac{Z_jk_ee^2}{2T}\right)^2
 \cot^2\frac{\theta_0}{2}.
\]
两成分单位面积核数
$n_{A,j}=w_j(\rho d)N_A/M_j$，
$w_{\rm Cu}=0.7,w_{\rm Zn}=0.3$，所以
\[
 \boxed{P_>=(\rho d)N_A
 \frac{\pi(k_ee^2)^2}{4T^2}\cot^2\frac{\theta_0}{2}
 \left(\frac{0.7Z_{\rm Cu}^2}{M_{\rm Cu}}+
 \frac{0.3Z_{\rm Zn}^2}{M_{\rm Zn}}\right)}.
\]
用 $Z_{\rm Cu}=29,Z_{\rm Zn}=30$、
$M_{\rm Cu}\simeq63.55\,\mathrm{g/mol}$、
$M_{\rm Zn}\simeq65.38\,\mathrm{g/mol}$，
\[
 \boxed{P_>\simeq1.4\times10^{-3}}.
\]
\excheck 黄铜两种元素应按\emph{质量份额/摩尔质量}
分别计算面核密度，不能直接按0.7、0.3平均原子序数。

\atomicproblem{16}{234}
求 $\alpha$ 粒子动能 $T=1.5\,\mathrm{MeV}$
在铀核（$Z=92$）库仑场散射偏角超过
$\theta_0=60^\circ$ 的有效截面。
\exsolution
积分卢瑟福微分截面，
\[
 \boxed{\sigma_{>\theta_0}=
 \pi\left(\frac{2Zk_ee^2}{2T}\right)^2
 \cot^2\frac{\theta_0}{2}}.
\]
由于 $k_ee^2=1.440\,\mathrm{MeV\,fm}$，
代入 $Z=92$、$\theta_0=60^\circ$，
得到
\[
 \boxed{\sigma_{>60^\circ}\simeq
 7.35\times10^{-26}\,\mathrm{m^2}
 =0.735\,\mathrm{kb}}.
\]
\excheck $1\,\mathrm{kb}=10^{-25}\,\mathrm{m^2}$；
本题需作角积分而非只报某个散射角的微分截面。

\atomicproblem{17}{234}
实验测得单个金核将单色 $\alpha$ 粒子散射在
$90^\circ$--$180^\circ$ 间的有效截面为
$\sigma=0.50\,\mathrm{kb}$。
求：(a) $\alpha$ 的动能；
(b) $\theta=60^\circ$ 时微分截面
$d\sigma/d\Omega$（$\mathrm{kb/sr}$）。
\exsolution
记 $C=2(79)k_ee^2$，
在后向半球
\[
 \sigma_{>90^\circ}=
 \pi\left(\frac C{2T}\right)^2.
\]
由此
\[
 \boxed{T=\frac C2\sqrt{\frac{\pi}{\sigma}}
 \simeq0.90\,\mathrm{MeV}}.
\]
而在 $\theta=60^\circ$，
\[
 \frac{d\sigma}{d\Omega}
 =\left(\frac C{4T}\right)^2
 \csc^4(30^\circ)
 =\frac{4\sigma_{>90^\circ}}{\pi}.
\]
所以
\[
 \boxed{\left.\frac{d\sigma}{d\Omega}\right|_{60^\circ}
 \simeq0.637\,\mathrm{kb/sr}}.
\]
\excheck 角积分量 $\sigma$ 与微分量
$d\sigma/d\Omega$ 的物理单位不同，不能混用。

\atomicproblem{18}{234}
由经典辐射功率公式估计：振动角频率
$\omega=5.0\times10^{15}\,\mathrm{s^{-1}}$
的电子近似简谐振动，其机械能减少为初值的
$1/\eta$（$\eta=10$）需多久？
\exsolution
电子加速度为
$a(t)=-\omega^2x_0\cos\omega t$。
SI 制拉莫尔辐射功率
$P(t)=e^2a^2/(6\pi\varepsilon_0c^3)$。
对一个振动周期取平均，
$\langle a^2\rangle=\omega^4x_0^2/2$；
振动机械能 $E=\tfrac12m_e\omega^2x_0^2$。
所以
\[
 \frac{dE}{dt}=-\frac{e^2\omega^2}
 {6\pi\varepsilon_0m_ec^3}E
 =-\Gamma E,\qquad
 E(t)=E(0)e^{-\Gamma t}.
\]
于是
\[
 \boxed{t=\frac{6\pi\varepsilon_0m_ec^3}
 {e^2\omega^2}\ln\eta
 \simeq1.5\times10^{-8}\,\mathrm s=15\,\mathrm{ns}}.
\]
\excheck 这是小损耗、准简谐振动的周期平均模型，
并已将原书高斯单位的辐射公式转换为 SI。
% === END 34-atomic-scattering-001-018.tex ===

% === BEGIN 34-atomic-scattering-019-032.tex blob 09661bcff99c46aebacb17ead9e95b132b4a5fd0 ===
% Irodov Part VI §6.1 problems 19--32.
% Original source printed pp.235--236, PDF pp.244--245.
% Distinguish principal quantum number n from ordinal line number:
% Balmer line k corresponds to upper n=k+2. This matters for #29.
\atomicproblem{19}{235}
用前题的经典辐射公式估计：一个电子沿半径
$r_0=50\,\mathrm{pm}$ 的氢原子圆轨道运动，
因连续辐射损失能量而落到核上需多久？
假定加速度始终指向核中心。
\exsolution
记库仑常数 $k_e=1/(4\pi\varepsilon_0)$。
电子圆轨道有
\[
 m_ev^2/r=k_ee^2/r^2,\qquad
 E(r)=\tfrac12m_ev^2-k_ee^2/r
 =-\frac{k_ee^2}{2r},
\]
并有 $a=k_ee^2/(m_er^2)$。
SI 制拉莫尔辐射功率
$P=e^2a^2/(6\pi\varepsilon_0c^3)
=2k_e^3e^6/(3m_e^2c^3r^4)$。
能量损失
$(dE/dr)(dr/dt)=-P$，
得到
\[
 \frac{dr}{dt}=-\frac{4k_e^2e^4}{3m_e^2c^3r^2}.
\]
从 $r_0$ 积分到0：
\[
 \boxed{t=\frac{m_e^2c^3r_0^3}{4k_e^2e^4}
 \simeq1.3\times10^{-11}\,\mathrm s=13\,\mathrm{ps}}.
\]
\excheck 这是古典模型的原子坍缩时间，
正是其无法描述真实稳定原子的经典矛盾。

\atomicproblem{20}{235}
证明：类氢离子电子从第 $n+1$ 轨道跃迁到第
$n$ 轨道时，辐射角频率 $\omega$
满足 $\Omega_n>\omega>\Omega_{n+1}$，
其中 $\Omega_j$ 是轨道上的电子公转角频率。
并证明大 $n$ 时 $\omega/\Omega_n\to1$。
\exsolution
玻尔能级写作 $E_n=-A/n^2$，$A>0$；
其第 $n$ 轨道公转角频率为
$\Omega_n=2A/(\hbar n^3)$。
相邻跃迁的光子角频率为
\[
 \omega=\frac{A}{\hbar}
 \left(\frac1{n^2}-\frac1{(n+1)^2}\right).
\]
对函数 $f(x)=-A/x^2$ 使用中值定理，
存在 $\xi\in(n,n+1)$ 使
$E_{n+1}-E_n=f'(\xi)=2A/\xi^3$，
因而
\[
 \boxed{\Omega_n>\omega>\Omega_{n+1}}.
\]
又 $\xi/n\to1$，得到
\[
 \boxed{\lim_{n\to\infty}
 \frac\omega{\Omega_n}=1}.
\]
\excheck 结论体现玻尔对应原理；
当 $n$ 大时相邻能级跃迁频率趋近经典公转频率。

\atomicproblem{21}{235}
质量为 $m$ 的粒子在中心势
$U(r)=kr^2/2$（$k>0$）中作圆周运动，
应用玻尔角动量量子化求允许的半径和能量。
\exsolution
圆周向心力条件
$mv^2/r=kr$，故公转角频率
$\Omega=v/r=\sqrt{k/m}$。
玻尔量子化
$L=mvr=mr^2\Omega=n\hbar$，
$n=1,2,\dots$，所以
\[
 \boxed{r_n=\sqrt{\frac{n\hbar}{m\Omega}}
 =\sqrt{\frac{n\hbar}{\sqrt{mk}}}}.
\]
机械能
$E_n=\tfrac12mv^2+\tfrac12kr_n^2
=kr_n^2$，于是
\[
 \boxed{E_n=n\hbar\Omega=
 n\hbar\sqrt{\frac km}}.
\]
\excheck 这是\emph{玻尔圆轨道}的模型，
不是一维量子谐振子精确本征值
$(n+\tfrac12)\hbar\Omega$。

\atomicproblem{22}{235}
对氢原子 H 和单电离氦离子 He$^+$，
求：(a) 第一玻尔轨道半径及电子速度；
(b) 基态动能与结合能；
(c) 电离电势、第一次激发电势、
共振谱线 $2\to1$ 波长。
\exsolution
固定核近似、核电荷数 $Z=1$ 或2，
\[
 r_{1,Z}=\frac{a_0}{Z},\quad
 v_{1,Z}=Z\alpha_{\rm fs}c,\quad
 K_{1,Z}=E_{b,Z}=13.606Z^2\,\mathrm{eV}.
\]
电离电势的数值（V）等于结合能的数值（eV）；
第一次激发所需能量为
$(1-1/4)E_b=3E_b/4$，
跃迁光子波长为 $\lambda_{21}=hc/(3E_b/4)$。
结果为
\[
\begin{array}{c|cc}
 &\mathrm H&\mathrm{He}^{+}\\\hline
 r_1\,(\mathrm{pm})&52.9&26.5\\
 v_1\,(10^6\mathrm{m/s})&2.18&4.36\\
 K_1=E_b\,(\mathrm{eV})&13.6&54.4\\
 U_{\mathrm{ion}}\,(\mathrm V)&13.6&54.4\\
 U_{\mathrm{exc},1}\,(\mathrm V)&10.2&40.8\\
 \lambda_{21}\,(\mathrm{nm})&121.6&30.4
\end{array}
\]
\excheck 原子 H 与类氢离子 He$^+$ 都只有一个电子，
同 $n$ 时 $r\propto Z^{-1}$、
$E\propto Z^2$、$\lambda\propto Z^{-2}$。

\atomicproblem{23}{235}
求单电离氦 He$^+$ 中电子沿第二玻尔轨道
运动的公转角频率。
\exsolution
类氢玻尔轨道半径
$r_n=a_0n^2/Z$，速度
$v_n=Z\alpha_{\rm fs}c/n$，
因此
\[
 \Omega_n=\frac{v_n}{r_n}
 =\frac{Z^2\alpha_{\rm fs}c}{a_0n^3}.
\]
对 He$^+$ 取 $Z=2,n=2$，
\[
 \boxed{\Omega_2\simeq2.07\times10^{16}\,\mathrm{s^{-1}}}.
\]
\excheck 求的是\emph{角频率}，
不是公转每秒圈数 $\Omega/(2\pi)$。

\atomicproblem{24}{235}
对类氢离子电子沿第 $n$ 个玻尔圆轨道运动，
求其轨道磁矩 $\boldsymbol\mu_n$、
轨道磁矩与轨道角动量之比；
给出第一轨道磁矩的模。
\exsolution
电子电荷为 $-e$，轨道等效电流
$I=-ev/(2\pi r)$，环路面积 $\pi r^2$，
故
\[
 \boldsymbol\mu=-\frac{evr}{2}\,\hat{\boldsymbol L}
 =-\frac e{2m_e}\boldsymbol L.
\]
按玻尔量子化 $|\boldsymbol L|=n\hbar$，
\[
 \boxed{\boldsymbol\mu_n=
 -n\frac{e\hbar}{2m_e}\,\hat{\boldsymbol L}},
 \qquad
 \boxed{\frac{|\boldsymbol\mu_n|}{|\boldsymbol L|}
 =\frac e{2m_e}},
\]
\[
 \boxed{|\boldsymbol\mu_1|=\mu_B=
 \frac{e\hbar}{2m_e}
 \simeq9.274\times10^{-24}\,\mathrm{J/T}}.
\]
\excheck 电子带负电，轨道磁矩矢量与
角动量反向；磁矩模是正值。

\atomicproblem{25}{235}
在玻尔模型中估算氢原子基态电子运动在
轨道中心产生的磁感应强度。
\exsolution
将电子轨道视为半径
$a_0=5.29\times10^{-11}\,\mathrm m$ 的圆电流环，
电流模为 $I=ev/(2\pi a_0)$，
环中心磁场为
\[
 B=\frac{\mu_0 I}{2a_0}
 =\frac{\mu_0 ev}{4\pi a_0^2}.
\]
用 $v=\alpha_{\rm fs}c\simeq2.19\times10^6\,\mathrm{m/s}$，
\[
 \boxed{B\simeq12.5\,\mathrm T
 =1.25\times10^5\,\mathrm G}.
\]
\excheck 这里估算的是经典玻尔\emph{轨道电流}磁场，
真实基态电子没有经典意义上的绕核轨道；
不应把该数值说成精确量子自旋/轨道局域磁场。

\atomicproblem{26}{235}
计算氢原子莱曼系、巴耳末系、帕邢系
在波长轴上覆盖的区间，并标出可见光区域。
\exsolution
里德伯公式
\[
 \frac1\lambda=R_H\left(\frac1{n_1^2}-\frac1{n_2^2}\right),
 \qquad n_2>n_1.
\]
对同一系，最长波长来自 $n_2=n_1+1$，
最短波长（系限）来自 $n_2\to\infty$，
故
\[
 \lambda_{\max}=
 \frac1{R_H[n_1^{-2}-(n_1+1)^{-2}]},
 \quad
 \lambda_{\min}=\frac{n_1^2}{R_H}.
\]
取近似 $R_H=1.097\times10^7\,\mathrm{m^{-1}}$：
\[
 \boxed{\begin{array}{c|cc}
 &\lambda_{\min}&\lambda_{\max}\\\hline
 \text{莱曼 }(n_1=1)&91.2\,\mathrm{nm}&121.6\,\mathrm{nm}\\
 \text{巴耳末 }(n_1=2)&364.6\,\mathrm{nm}&656.3\,\mathrm{nm}\\
 \text{帕邢 }(n_1=3)&820.4\,\mathrm{nm}&1875\,\mathrm{nm}
 \end{array}}
\]
\excheck 可见区约为 $380$--$760\,\mathrm{nm}$，
其中的氢谱线主要属于巴耳末系；
该表可直接据端点在波长轴绘制三个连续区间。

\atomicproblem{27}{235}
氢原子某条谱线的波数，
等于巴耳末系波长
$486.1\,\mathrm{nm}$ 与 $410.2\,\mathrm{nm}$
两条谱线波数之差。
求该未知谱线属于哪一系及其波长。
\exsolution
巴耳末 $486.1\,\mathrm{nm}$ 为 $4\to2$，
$410.2\,\mathrm{nm}$ 为 $6\to2$。
其波数之\emph{差}为
\[
 \left[\frac1{410.2\,\mathrm{nm}}-
 \frac1{486.1\,\mathrm{nm}}\right]
 =R_H\left(\frac1{4^2}-\frac1{6^2}\right).
\]
它对应从上能级6跃迁到下能级4，
即\emph{布拉开特系（Brackett）}的 $6\to4$ 谱线。
于是
\[
 \boxed{\lambda=\frac1{R_H(1/16-1/36)}
 \simeq2.63\,\mu\mathrm m}.
\]
\excheck 必须取两谱线\emph{波数}之差，
而不是直接拿两个波长相减。

\atomicproblem{28}{235}
对氢原子求：(a) 巴耳末系前三条谱线波长；
(b) 能分辨该系列前20条谱线的光谱仪，
至少应有多大的分辨本领 $\lambda/\delta\lambda$？
\exsolution
巴耳末系下能级为 $n_1=2$，
第 $k$ 条线对应上能级 $n_2=k+2$。
由里德伯公式，
\[
 \boxed{\lambda_1\simeq656\,\mathrm{nm},\quad
 \lambda_2\simeq486\,\mathrm{nm},\quad
 \lambda_3\simeq434\,\mathrm{nm}}.
\]
高序线间距不断缩小，最严格条件是区分
第20条（$22\to2$）与第21条（$23\to2$）。
故要求
\[
 \mathcal R\ge
 \frac{\lambda_{22}}{\lambda_{22}-\lambda_{23}},
 \quad
 \lambda_n=\frac1{R_H(1/4-1/n^2)}.
\]
直接计算
\[
 \boxed{\mathcal R_{\min}\simeq1.41\times10^3
 \approx1.5\times10^3}.
\]
\excheck 应比较最密集的一端；
不能只用第一条和第二条的大间隔估计。

\atomicproblem{29}{235}
宽度 $\ell=6.6\,\mathrm{mm}$ 的光栅
受到正入射的氢原子辐射。
巴耳末系的\emph{第50条谱线}在衍射角 $\theta$
附近恰达到瑞利分辨判据。
求 $\theta$，并对“第50条”序号作严格辨析。
\exsolution
令光栅实际参与干涉的缝数为 $N$，
衍射级次为 $m$。
光栅的瑞利分辨本领为
$\mathcal R=mN$，而正入射光栅
$m\lambda=d\sin\theta$、$N\simeq\ell/d$，
所以
\[
 \boxed{\sin\theta=
 \frac{\lambda^2}{\ell\,\delta\lambda}}.
\]
上式中的分子分母约去后实质为
$\sin\theta=\lambda^2/(\ell\delta\lambda)$。
严格按\emph{第50条}定义，
上能级 $n=52$，邻线是 $n=53$。
用里德伯公式
$\lambda_{52}\simeq365.047\,\mathrm{nm}$、
$\delta\lambda\simeq0.02021\,\mathrm{nm}$，
则
\[
 \boxed{\sin\theta\simeq0.999,\quad
 \theta\simeq87^\circ}.
\]
\textbf{原书答案约 $60^\circ$ 的来源：}
若把“第50条”\emph{直接当作上能级 $n=50$}，
则相邻谱线 $50\to2$、$51\to2$
给 $\sin\theta\simeq0.889$、$\theta\simeq63^\circ$，
即书末近似 $60^\circ$。
\excheck 第50条谱线与 $n=50$ 上能级
原则上差2；本题的原书题号解释需由独立学术复核决定，
不能悄然抹去两种编号造成的差异。

\atomicproblem{30}{236}
哪种元素的类氢谱线波长等于氢原子对应谱线的四分之一？
\exsolution
类氢离子（忽略约化质量细修正）的能级
$E_n\propto-Z^2/n^2$，
对应两能级间辐射的光子能量
$\Delta E\propto Z^2$。
因此对应波长
$\lambda_Z=\lambda_H/Z^2$。
由 $\lambda_Z=\lambda_H/4$ 得 $Z=2$，
所以
\[
 \boxed{\mathrm{He}^{+}\ \text{（单电离氦离子）}}.
\]
\excheck 普通氦原子有两个电子，
不是类氢单电子体系；题中对应的是 He$^+$。

\atomicproblem{31}{236}
氢原子被激发至主量子数 $n$ 的能级。
假定允许其在各个更低能级间辐射跃迁，
最多有多少条不同频率的谱线？
\exsolution
可参与级联跃迁的能级集合为
$1,2,\dots,n$。
每一对不同能级 $(i,j)$，$1\le i<j\le n$，
对应一次可能的向下跃迁，
故谱线种数为二元组合数
\[
 \boxed{N=\binom n2=\frac{n(n-1)}{2}}.
\]
\excheck 本题数的是\emph{整组被激发原子}
可能发出的谱线，不是单个原子一次级联必然发出的线数。

\atomicproblem{32}{236}
求氢原子\emph{吸收}谱在
$94.5$--$130.0\,\mathrm{nm}$ 区间内有哪些线。
初始原子均在基态。
\exsolution
基态吸收只能有 $1\to n$，$n\ge2$，
属于莱曼系，满足
\[
 \lambda_{1n}=
 \frac1{R_H(1-1/n^2)}.
\]
依次为
\[
 \boxed{\lambda_{12}\simeq121.6\,\mathrm{nm}},\quad
 \boxed{\lambda_{13}\simeq102.6\,\mathrm{nm}},\quad
 \boxed{\lambda_{14}\simeq97.3\,\mathrm{nm}},\quad
 \boxed{\lambda_{15}\simeq95.0\,\mathrm{nm}}.
\]
下一条 $1\to6$ 的波长约 $93.8\,\mathrm{nm}$，
已小于题给下界 $94.5\,\mathrm{nm}$。
\excheck 按底本题面区间应有\emph{四条}，
而底本书末只列前三条
$97.3,102.6,121.6\,\mathrm{nm}$，
漏掉仍在区间内的 $95.0\,\mathrm{nm}$；
已明确列为需校勘的答案缺项。
% === END 34-atomic-scattering-019-032.tex ===

% === BEGIN 34-atomic-scattering-033-048.tex blob 26103729fd44a2e0cd4bcce44248fe70e51a9801 ===
% Irodov Part VI §6.1 problems 33--48 (end of §6.1).
% Source: Arihant PDF pp.245--246 (printed pp.236--237),
% Answers pp.349 (printed pp.340). Source scans establish
% #36=5.18e15 s^-1, #39=18.0 nm, #44=0.20 nm.
\atomicproblem{33}{236}
He$^+$ 离子从某一激发态连续发射两光子，
其波长分别为 $108.5$ 与 $30.4\,\mathrm{nm}$，
最终达到基态。求初始主量子数 $n$。
\exsolution
He$^+$ 基态结合能约 $E_b=54.4\,\mathrm{eV}$。
由跃迁能量
\[
 \Delta E=E_b\left(\frac1{n_f^2}-\frac1{n_i^2}\right)
 =\frac{hc}{\lambda}
\]
知 $30.4\,\mathrm{nm}$ 对应约
$40.8\,\mathrm{eV}$，
即 $2\to1$ 跃迁。
$108.5\,\mathrm{nm}$ 对应约
$11.43\,\mathrm{eV}$，因此前一跃迁满足
\[
 11.43=54.4\left(\frac14-\frac1{n^2}\right).
\]
解得 $1/n^2\simeq0.040$，
\[
 \boxed{n=5},\quad
 \text{跃迁顺序 }5\to2\to1.
\]
\excheck 两光子能量之和
$11.43+40.8\simeq52.2\,\mathrm{eV}$，
等于 $54.4(1-1/25)\,\mathrm{eV}$。

\atomicproblem{34}{236}
已知 He$^+$ 的巴耳末系和莱曼系第一条
（各自最长波长）谱线波长之差为
$\Delta\lambda=133.7\,\mathrm{nm}$，
求以\emph{角频率}记号定义的里德伯常数
$\mathcal R$。
\exsolution
类氢离子跃迁角频率写作
$\omega=\mathcal R Z^2(n_1^{-2}-n_2^{-2})$，
对 He$^+$ 取 $Z=2$。
第一条巴耳末为 $3\to2$，
第一条莱曼为 $2\to1$，因此
\[
 \lambda_B=\frac{2\pi c}{\mathcal RZ^2(1/4-1/9)}
 =\frac{72\pi c}{5\mathcal RZ^2},
\quad
 \lambda_L=\frac{2\pi c}{\mathcal RZ^2(1-1/4)}
 =\frac{8\pi c}{3\mathcal RZ^2}.
\]
相减得
\[
 \Delta\lambda=\frac{176\pi c}{15\mathcal RZ^2},
 \qquad
 \boxed{\mathcal R=
 \frac{176\pi c}{15Z^2\Delta\lambda}
 \simeq2.07\times10^{16}\,\mathrm{s^{-1}}}.
\]
\excheck 这里 $\mathcal R$ 是\emph{角频率常数}，
与常用波数里德伯常数
$R_\infty\simeq1.097\times10^7\,\mathrm{m^{-1}}$
相差因子 $2\pi c$。

\atomicproblem{35}{236}
某单电子类氢离子巴耳末系与莱曼系
第一谱线的波长差为 $59.3\,\mathrm{nm}$。
确定该离子。
\exsolution
由上一题可知，对于给定跃迁组合，
波长差按核电荷 $Z$ 缩放为
$\Delta\lambda_Z\propto1/Z^2$。
He$^+$（$Z=2$）时差为 $133.7\,\mathrm{nm}$，
于是
\[
 \left(\frac Z2\right)^2=
 \frac{133.7}{59.3},\quad
 Z\simeq3.00.
\]
因此为
\[
 \boxed{\mathrm{Li}^{2+}
 \quad\text{（双电离锂，核电荷数3）}}.
\]
\excheck 中性锂原子含3电子，不属于类氢单电子体系。

\atomicproblem{36}{236}
He$^+$ 的某一谱系，其最大与最小
跃迁角频率之差为
$\Delta\omega=5.18\times10^{15}\,\mathrm{s^{-1}}$。
求这个谱系的第一条（最长波长）谱线波长。
\exsolution
设谱系下能级为 $n$。
系限角频率为
$\omega_\infty=\mathcal R Z^2/n^2$，
第一线来自 $n+1\to n$，
$\omega_1=\mathcal RZ^2[n^{-2}-(n+1)^{-2}]$。
故间隔
\[
 \Delta\omega=\omega_\infty-\omega_1
 =\frac{\mathcal R Z^2}{(n+1)^2}.
\]
对 $Z=2$、$\mathcal R=2.07\times10^{16}\,\mathrm{s^{-1}}$，
得到 $(n+1)^2\simeq16$，所以 $n=3$。
第一线为 $4\to3$，
\[
 \boxed{\lambda_{43}=
 \frac{2\pi c}{\mathcal R\,4(1/9-1/16)}
 \simeq0.47\,\mu\mathrm m}.
\]
\excheck 题中是两端\emph{角频率}之差，
不可用普通频率常数遗漏 $2\pi$。

\atomicproblem{37}{236}
某类氢离子巴耳末系第三条谱线波长
为 $108.5\,\mathrm{nm}$。
求该离子基态电子结合能。
\exsolution
巴耳末第三线为 $5\to2$，
对应光子能量
\[
 \frac{hc}{\lambda}
 =E_b\left(\frac14-\frac1{25}\right)
 =\frac{21}{100}E_b.
\]
所以
\[
 \boxed{E_b=\frac{100}{21}
 \frac{hc}{108.5\,\mathrm{nm}}
 \simeq54.4\,\mathrm{eV}}.
\]
该值对应 He$^+$。
\excheck 不同谱线序数须区别于主量子数；
第三条巴耳末线的上能级为5。

\atomicproblem{38}{236}
中性氦原子基态第一个电子的结合能
（第一电离能）为
$E_0=24.6\,\mathrm{eV}$。
求移去氦原子全部两个电子所需总能量。
\exsolution
先移去第一个电子，消耗
$24.6\,\mathrm{eV}$；
剩余 He$^+$ 是核电荷为2的类氢离子，
其基态第二电离能为
$4\times13.6=54.4\,\mathrm{eV}$。
所以
\[
 \boxed{E_{\rm total}=24.6+54.4
 =79.0\,\mathrm{eV}}.
\]
\excheck 不能把第一电离能简单乘2，
因为第一次电离前后电子间排斥的影响不同。

\atomicproblem{39}{236}
波长 $\lambda=18.0\,\mathrm{nm}$ 的辐射
使起初静止的 He$^+$ 离子基态发生光电离。
求光电子最大速度。
\exsolution
单个入射光子能量
$hc/\lambda\simeq68.88\,\mathrm{eV}$，
He$^+$ 的结合能 $E_b=54.4\,\mathrm{eV}$。
忽略重核反冲，
电子动能
$K=hc/\lambda-E_b\simeq14.48\,\mathrm{eV}$。
故
\[
 \boxed{v_{\max}=
 \sqrt{\frac{2(hc/\lambda-E_b)}{m_e}}
 \simeq2.26\times10^6\,\mathrm{m/s}}.
\]
\excheck 原扫描题面为 \textbf{$18.0\,\mathrm{nm}$}，
不是 PDF 抽取文本易误识别的 $180\,\mathrm{nm}$；
后者光子能量不足以电离基态 He$^+$。

\atomicproblem{40}{236}
基态氢原子 A 以初始动能 $T$ 正面撞击
另一个起初静止的基态氢原子 B。
求使其中一个原子至少激发到
$n=2$、进而能辐射光子的最小 $T$。
\exsolution
使一个氢原子从基态激发到 $n=2$，
需要内部激发能
$\Delta=13.6(1-1/4)=10.2\,\mathrm{eV}$。
两个原子质量相等 $M$，
实验室系中可转化为内部激发的质心相对运动能量
为
\[
 T_{\rm cm}=\frac12\mu v^2
 =\frac12\frac M2\,v^2
 =\frac12T,\qquad T=\frac12Mv^2.
\]
所以临界条件 $T/2=\Delta$，
\[
 \boxed{T_{\min}=2\Delta\simeq20.4\,\mathrm{eV}}.
\]
\excheck 不能仅取 $10.2\,\mathrm{eV}$，
因为碰后质心必须继续运动，
入射动能不可能全部转为内部激发。

\atomicproblem{41}{236}
起初静止的氢原子放出第一条莱曼线光子。
求氢原子获得的反冲速度。
\exsolution
第一莱曼线来自 $2\to1$，
光子能量
$\varepsilon\simeq10.2\,\mathrm{eV}$，
动量 $p_\gamma=\varepsilon/c$。
若原子起初静止，由总动量守恒，
原子获得等大反向动量，故
\[
 \boxed{v_{\mathrm H}=
 \frac{\varepsilon}{M_{\mathrm H}c}
 \simeq3.25\,\mathrm{m/s}}.
\]
\excheck 原子质量使用氢\emph{原子}质量；
光子的质量不是 $E/c^2$ 再去用非相对论动量公式求速度。

\atomicproblem{42}{236}
承上题，求实际发出光子的能量
与不计反冲时的 $2\to1$ 内能差之间的相对差，
以百分数表示。
\exsolution
跃迁内能差记为 $\Delta E$。
发射后原子后退动能
$K_{\rm recoil}=p_\gamma^2/(2M)
=E_\gamma^2/(2Mc^2)$。
能量守恒
$\Delta E=E_\gamma+E_\gamma^2/(2Mc^2)$；
因 $\Delta E\ll Mc^2$，
一阶近似
\[
 \frac{\Delta E-E_\gamma}{\Delta E}
 \simeq\frac{\Delta E}{2Mc^2}.
\]
以 $\Delta E=10.2\,\mathrm{eV}$、
$Mc^2\simeq938.8\,\mathrm{MeV}$，
\[
 \boxed{\text{相对差}\simeq
 5.43\times10^{-9}
 =5.43\times10^{-7}\%\approx0.55\times10^{-6}\%}.
\]
\excheck 这是一个极小的\emph{红移}校正；
将比值改写百分数时必须乘100。

\atomicproblem{43}{236}
起初静止的 He$^+$ 发出第一莱曼线光子，
该光子随后电离起初静止的基态氢原子。
求被释放光电子的速度，忽略重核反冲。
\exsolution
He$^+$ 的 $2\to1$ 光子能量为
$4(10.2)=40.8\,\mathrm{eV}$。
氢基态电离能为
$13.6\,\mathrm{eV}$。
由能量守恒，
光电子获得动能
$T_e=40.8-13.6=27.2\,\mathrm{eV}$，
因而
\[
 \boxed{v_e=\sqrt{\frac{2T_e}{m_e}}
 \simeq3.09\times10^6\,\mathrm{m/s}}.
\]
\excheck 光子具有离散跃迁能，
不能将 He$^+$ 的 $40.8\,\mathrm{eV}$ 误作光电子动能；
必须扣除氢的电离能。

\atomicproblem{44}{236}
沿与原子速度成 $\theta=45^\circ$
的方向观察激发氢原子的第一莱曼线，
发现波长位移 $\Delta\lambda=0.20\,\mathrm{nm}$。
估算氢原子速率。
\exsolution
在 $v\ll c$ 时，相对论多普勒公式的一阶形式为
\[
 \frac{|\Delta\lambda|}{\lambda_0}
 \simeq\frac{v|\cos\theta|}{c}.
\]
第一莱曼线
$\lambda_0\simeq121.6\,\mathrm{nm}$，
因此
\[
 \boxed{v\simeq
 \frac{c\,|\Delta\lambda|}
 {\lambda_0|\cos45^\circ|}
 \simeq6.98\times10^5\,\mathrm{m/s}}.
\]
\excheck 位移的符号由接近或远离决定；
题给只问速度大小，故取绝对值。

\atomicproblem{45}{236--237}
应用玻尔--索末菲作用量量子化
$\oint p\,dq=2\pi n\hbar$，
求粒子在：(a) 长度 $\ell$ 的无限深一维方阱；
(b) 半径 $r$ 的圆周；
(c) $U(x)=\alpha x^2/2$ 的一维谐振势；
(d) $U(r)=-\alpha/r$ 的中心引力型势作圆周运动时，
允许的能量。
\exsolution
(a) 粒子在方阱两侧来回反射，一个周期走过 $2\ell$，
故 $2p\ell=2\pi n\hbar$，
\[
 \boxed{E_n=\frac{n^2\pi^2\hbar^2}{2m\ell^2}}.
\]
(b) 沿圆周 $\oint p\,dq=2\pi rp=2\pi L$，
所以 $L=n\hbar$，有
\[
 \boxed{E_n=\frac{n^2\hbar^2}{2mr^2}}.
\]
(c) 振子频率 $\Omega=\sqrt{\alpha/m}$，
经典一周作用量
$\oint p\,dx=2\pi E/\Omega$，于是
\[
 \boxed{E_n=n\hbar\Omega=n\hbar\sqrt{\alpha/m}}.
\]
(d) 圆周条件 $mv^2/r=\alpha/r^2$，
加上 $L=mvr=n\hbar$，
解出
$r=n^2\hbar^2/(m\alpha)$，
能量为
\[
 \boxed{E_n=-\frac{m\alpha^2}{2n^2\hbar^2}}.
\]
\excheck 这是\emph{早期作用量量子化}推导；
(c) 缺少精确量子谐振子的零点能
$\hbar\Omega/2$，不应把这两种理论混为一谈。

\atomicproblem{46}{237}
考虑氢原子质子的运动，
求基态电子结合能及里德伯常数的
约化质量修正，并估计忽略核运动的相对误差。
\exsolution
电子与质子的相对运动等价于约化质量
\[
 \mu=\frac{m_em_p}{m_e+m_p}
\]
的单粒子库仑问题。
因玻尔结合能与里德伯常数都正比于质量，
\[
 \boxed{E_b=\frac{\mu(k_ee^2)^2}{2\hbar^2}},\qquad
 \boxed{\mathcal R=
 \frac{\mu(k_ee^2)^2}{2\hbar^3}},
\]
其中 $\mathcal R$ 是角频率形式，
若使用波数里德伯常数需除以 $2\pi c$。
忽略质子运动等于用 $m_e$ 代替 $\mu$，
相对高估量为
\[
 \frac{m_e-\mu}{\mu}=\frac{m_e}{m_p}
 \simeq5.45\times10^{-4}.
\]
故
\[
 \boxed{\text{相对误差}\simeq0.0545\%}.
\]
\excheck 对真实氢原子这一修正很小，
但分析同位素光谱线位移时不可遗漏。

\atomicproblem{47}{237}
对普通氢 H 和重氢 D，
计算：(a) 基态电子结合能之差；
(b) 第一莱曼线波长之差。
\exsolution
两体系电子质量相同，核质量分别为
$m_p$、$m_d$，约化质量
$\mu_A=m_em_A/(m_e+m_A)$。
结合能及跃迁频率都与 $\mu$ 成正比，
故
\[
 \boxed{E_{b,D}-E_{b,H}
 =13.6057\,\mathrm{eV}\,
 \frac{\mu_D-\mu_H}{m_e}
 \simeq3.7\,\mathrm{meV}}.
\]
第一莱曼线波长
$\lambda_{21,A}\propto1/\mu_A$，
因此
\[
 \boxed{\lambda_{21,H}-\lambda_{21,D}
 \simeq33\,\mathrm{pm}}.
\]
\excheck 重氢核较重，约化质量更接近 $m_e$，
因此 D 的结合能稍大、莱曼波长稍短。

\atomicproblem{48}{237}
估计以下两个类氢系统基态时
两粒子间距、结合能和第一莱曼线波长：
(a) 质子与带电量 $-e$、质量 $207m_e$ 的负介子
组成的介子氢；
(b) 电子与正电子组成的正电子素。
\exsolution
对电荷模为 $e$ 的两粒子库仑体系，
只需把电子质量换成约化质量
$\mu=m_1m_2/(m_1+m_2)$。
基态类玻尔半径、结合能、第一莱曼线依次为
\[
 a=\frac{4\pi\varepsilon_0\hbar^2}{\mu e^2},
 \quad E_b=\frac{\mu(k_ee^2)^2}{2\hbar^2},
 \quad\lambda_{21}=\frac{hc}{3E_b/4}.
\]
(a) 对 $\mu_{\mu p}=
(207m_e)m_p/(207m_e+m_p)\simeq186m_e$，
\[
 \boxed{a\simeq0.285\,\mathrm{pm},\quad
 E_b\simeq2.53\,\mathrm{keV},\quad
 \lambda_{21}\simeq0.65\,\mathrm{nm}}.
\]
(b) 正电子素约化质量恰为 $\mu=m_e/2$，
所以
\[
 \boxed{a\simeq106\,\mathrm{pm},\quad
 E_b\simeq6.80\,\mathrm{eV},\quad
 \lambda_{21}\simeq243\,\mathrm{nm}}.
\]
\excheck 第(a)问不能直接用介子质量 $207m_e$
代替约化质量，质子质量有限会产生明显修正。
% === END 34-atomic-scattering-033-048.tex ===
